% math0 manuscript. Build with: latexmk -pdf paper.tex
\documentclass[11pt]{amsart}
\usepackage[margin=1.1in]{geometry}
\usepackage{amsmath,amssymb,amsthm,mathtools}
\usepackage{url}
\usepackage{hyperref}
\usepackage{xcolor}
\usepackage{textcomp}
\usepackage{listings}
\usepackage{newunicodechar}

\lstdefinelanguage{lean}{
  keywords={theorem,lemma,def,abbrev,example,fun,by,have,show,from,let,in,if,then,else,match,with,
    namespace,end,open,import,section,variable,instance,structure,where,calc,obtain,intro,exact,
    refine,rcases,rintro,simp,omega,decide,norm_num,linarith,nlinarith,positivity,induction,cases,
    private,rw},
  sensitive=true, comment=[l]{--}, morecomment=[s]{/-}{-/}, string=[b]"}
\lstset{language=lean, basicstyle=\ttfamily\footnotesize, columns=fullflexible, keepspaces=true, breaklines=true,
  keywordstyle=\color{blue!60!black}\bfseries, commentstyle=\color{gray}, frame=single,
  extendedchars=true, inputencoding=utf8, upquote=true,
  literate={ℕ}{{$\mathbb{N}$}}1 {∀}{{$\forall$}}1 {↔}{{$\leftrightarrow$}}1
    {¬}{{$\neg$}}1 {∧}{{$\wedge$}}1 {∨}{{$\vee$}}1 {→}{{$\to$}}1 {≤}{{$\le$}}1
    {⟨}{{$\langle$}}1 {⟩}{{$\rangle$}}1}
% Unicode commonly found in Lean sources.
\newunicodechar{ℕ}{\ensuremath{\mathbb{N}}}
\newunicodechar{ℤ}{\ensuremath{\mathbb{Z}}}
\newunicodechar{ℚ}{\ensuremath{\mathbb{Q}}}
\newunicodechar{ℝ}{\ensuremath{\mathbb{R}}}
\newunicodechar{ℂ}{\ensuremath{\mathbb{C}}}
\newunicodechar{∀}{\ensuremath{\forall}}
\newunicodechar{∃}{\ensuremath{\exists}}
\newunicodechar{→}{\ensuremath{\to}}
\newunicodechar{←}{\ensuremath{\leftarrow}}
\newunicodechar{↔}{\ensuremath{\leftrightarrow}}
\newunicodechar{↦}{\ensuremath{\mapsto}}
\newunicodechar{¬}{\ensuremath{\neg}}
\newunicodechar{∧}{\ensuremath{\wedge}}
\newunicodechar{∨}{\ensuremath{\vee}}
\newunicodechar{≤}{\ensuremath{\le}}
\newunicodechar{≥}{\ensuremath{\ge}}
\newunicodechar{≠}{\ensuremath{\ne}}
\newunicodechar{∈}{\ensuremath{\in}}
\newunicodechar{∉}{\ensuremath{\notin}}
\newunicodechar{⊆}{\ensuremath{\subseteq}}
\newunicodechar{∩}{\ensuremath{\cap}}
\newunicodechar{∪}{\ensuremath{\cup}}
\newunicodechar{∑}{\ensuremath{\sum}}
\newunicodechar{∏}{\ensuremath{\prod}}
\newunicodechar{⟨}{\ensuremath{\langle}}
\newunicodechar{⟩}{\ensuremath{\rangle}}
\newunicodechar{·}{\ensuremath{\cdot}}
\newunicodechar{∘}{\ensuremath{\circ}}
\newunicodechar{⁻}{\ensuremath{^{-}}}
\newunicodechar{¹}{\ensuremath{^{1}}}
\newunicodechar{²}{\ensuremath{^{2}}}
\newunicodechar{₀}{\ensuremath{_{0}}}
\newunicodechar{₁}{\ensuremath{_{1}}}
\newunicodechar{₂}{\ensuremath{_{2}}}
\newunicodechar{α}{\ensuremath{\alpha}}
\newunicodechar{β}{\ensuremath{\beta}}
\newunicodechar{ε}{\ensuremath{\varepsilon}}
\newunicodechar{δ}{\ensuremath{\delta}}
\newunicodechar{π}{\ensuremath{\pi}}
\newunicodechar{σ}{\ensuremath{\sigma}}
\newunicodechar{φ}{\ensuremath{\varphi}}
\newunicodechar{∞}{\ensuremath{\infty}}
\newunicodechar{⌊}{\ensuremath{\lfloor}}
\newunicodechar{⌋}{\ensuremath{\rfloor}}
\newunicodechar{⌈}{\ensuremath{\lceil}}
\newunicodechar{⌉}{\ensuremath{\rceil}}
\newunicodechar{∣}{\ensuremath{\mid}}
\newunicodechar{≡}{\ensuremath{\equiv}}
\newunicodechar{×}{\ensuremath{\times}}
\newunicodechar{⧸}{\ensuremath{/}}
\newunicodechar{‖}{\ensuremath{\|}}
\newunicodechar{√}{\ensuremath{\surd}}
\newunicodechar{∅}{\ensuremath{\emptyset}}
\newunicodechar{⦃}{\ensuremath{\{\!|}}
\newunicodechar{⦄}{\ensuremath{|\!\}}}
\newunicodechar{𝓝}{\ensuremath{\mathcal{N}}}
\newunicodechar{≫}{\ensuremath{\gg}}
\newunicodechar{≪}{\ensuremath{\ll}}

\theoremstyle{plain}
\newtheorem{theorem}{Theorem}
\newtheorem{lemma}[theorem]{Lemma}
\newtheorem{proposition}[theorem]{Proposition}
\newtheorem{corollary}[theorem]{Corollary}
\theoremstyle{definition}
\newtheorem{definition}[theorem]{Definition}
\newtheorem{remark}[theorem]{Remark}

\emergencystretch=3em
% Lean identifiers in running text (breakable at dots and underscores).
\expandafter\def\expandafter\UrlBreaks\expandafter{\UrlBreaks\do\_\do\.}
\DeclareUrlCommand\decl{\urlstyle{tt}}
\newcommand{\N}{\mathbb{N}}
\newcommand{\Z}{\mathbb{Z}}
\DeclareMathOperator{\ord}{ord}

\title[Repunit residues and long period primes]{The repunits $R_1,\dots,R_n$ leave exactly $n-1$ distinct residues modulo $n$ if and only if $n$ is a long period prime}
\author{math0}
\address{Produced by Claude (an AI model developed by Anthropic) in the math0 project}
\date{October 18, 2026}

\begin{document}
\begin{abstract}
For $n\ge1$ let $a(n)$ be the number of distinct remainders of the repunits $R_k=(10^k-1)/9$, $1\le k\le n$, upon division by $n$; this is OEIS A335012.
We prove the conjecture recorded in that entry by Sen-Peng Eu (2020) that $a(n)=n-1$ if and only if $n$ is a long period prime, that is, $n$ is prime and the decimal expansion of $1/n$ has period $n-1$ (OEIS A006883, which includes $2$).
The proof is elementary: it uses the periodicity of repunits modulo $n$, the congruence $R_k\equiv k\pmod 3$, and Euler's theorem.
To our knowledge it is the first complete proof; the OEIS entry still lists the statement as a conjecture, and the case $\gcd(n,30)=1$ follows readily from a formula stated without proof by Robert Israel in the same entry.
The theorem is verified in Lean~4 with Mathlib against a statement that math0 formalized from the OEIS entry, whose definitions are checked by the kernel against the listed OEIS terms; this statement is not part of the formal-conjectures repository.
\end{abstract}
\maketitle

\section{Introduction}

For $k\ge0$ let $R_k=(10^k-1)/9$ be the $k$-th repunit, so $R_1=1$, $R_2=11$, $R_3=111$, and so on (OEIS A002275 \cite{A002275}).
The sequence A335012 of the OEIS \cite{A335012}, submitted by Sen-Peng Eu on May~19, 2020, has offset $1$ and is defined by
\[
  a(n)=\#\{\,R_k \bmod n \;:\; 1\le k\le n\,\},
\]
the ``number of different remainders when the first $n$ terms of $1, 11, 111, 1111, \dots$ are divided by $n$''.
For example $a(4)=2$, since $1,11,111,1111$ leave the remainders $1,3,3,3$ modulo $4$.
The sequence begins
\[
  1, 1, 3, 2, 1, 3, 6, 3, 9, 1, 2, 4, 6, 6, 3, 4, 16, 9, 18, 2, 6, 2, 22, \dots
\]
Since there are only $n$ residues modulo $n$, $a(n)\le n$.

A prime $p$ is a \emph{long period prime} if the decimal expansion of $1/p$ has period $p-1$ (OEIS A006883 \cite{A006883}).
By the convention stated in A006883, $2$ is a long period prime, because $1/2=0.5000\ldots$ is eventually periodic with period $1=2-1$; the related sequence A001913 \cite{A001913} of full reptend primes, the primes with primitive root $10$, omits $2$.
For a prime $p\notin\{2,5\}$ the period of $1/p$ equals the multiplicative order $\ord_p(10)$ of $10$ modulo $p$ \cite[Ch.~IX]{HW}, so $p$ is a long period prime exactly when $10$ is a primitive root modulo $p$ (see also \cite{MW}).
Throughout, for $n\ge1$, $\ord_n(10)$ denotes the least $m\ge1$ with $10^m\equiv1\pmod n$; it exists if and only if $\gcd(n,10)=1$, and when it does not exist an equation $\ord_n(10)=n-1$ is regarded as false (this agrees with Mathlib's convention, see Section~\ref{sec:statement}).
The primes $3$ and $5$ are not long period primes, since $1/3$ and $1/5$ have period $1$.
The long period primes begin $2,7,17,19,23,29,47,59,61,97,\dots$, and the values of $a$ listed above have $a(n)=n-1$ exactly at $n=2,7,17,19,23$.
The COMMENTS section of A335012 contains the following conjecture, which we prove.

\begin{theorem}\label{thm:main}
For every integer $n\ge1$, $a(n)=n-1$ if and only if $n$ is a long period prime, that is, $n$ is prime and either $n=2$ or $\ord_n(10)=n-1$.
\end{theorem}

We prove Theorem~\ref{thm:main} in Section~\ref{sec:proof} and verify it in Lean~4 (Section~\ref{sec:lean}).

\subsection*{Status and prior work}
The current revision of A335012 (revision \#31, March~12, 2026) still states the claim as a conjecture, and the entry has no extension line recording a proof or a counterexample.
The nearest earlier result is a comment by Robert Israel in the same entry, dated June~25, 2020: if $\gcd(n,30)=1$ then $a(n)$ equals A084680$(n)$ \cite{A084680}, the multiplicative order of $10$ modulo $n$.
The comment is stated without proof.
Combined with $\ord_n(10)\le\varphi(n)\le n-1$ for $n\ge2$, where equality throughout forces $n$ to be prime, it gives Theorem~\ref{thm:main} for $\gcd(n,30)=1$ (see Remark~\ref{rem:israel}, where we derive Israel's formula from the lemmas of Section~\ref{sec:proof}).
It says nothing about $n$ divisible by $2$, $3$ or $5$.
Multiples of $2$ and $5$ are easy to exclude; multiples of $3$ need the congruence $R_k\equiv k\pmod 3$.
The ingredients are standard (for instance, the pigeonhole argument showing that some repunit is divisible by any $n$ coprime to $10$, and $R_k\equiv k\pmod 3$), and the cross-referenced entries A006883 and A084680 do not mention a proof.
The math0 referees searched arXiv, MathOverflow, Mathematics Stack Exchange, the SeqFan mailing list and the general web, and found no earlier proof or counterexample.
The statement is not among the OEIS statements of the formal-conjectures repository \cite{FC} at the commit pinned by math0.
We therefore classify Theorem~\ref{thm:main} as a \emph{new result}, with the caveat that the argument is elementary and that its main case is essentially contained in Israel's comment; what is new is the treatment of the multiples of $2$, $3$ and $5$ and the complete proof of both directions.
The entry also records, in an unsigned comment, that $a(n)=n$ if and only if $n$ is a power of $3$; we do not address that statement.

\section{Statement}\label{sec:statement}

The COMMENTS section of A335012 \cite{A335012} contains, verbatim:
\begin{quote}\raggedright
{\small\texttt{Conjecture: a(n) = n-1 if and only if n is a long period prime (A006883), that is, n is a prime and the decimal expansion of 1/n has period n-1.}}
\end{quote}
The comment carries no separate signature, so we attribute it to the author of the entry, Sen-Peng Eu (May~19, 2020).
Theorem~\ref{thm:main} is this conjecture, read for every index $n\ge1$ of the sequence and with long period primes as defined in A006883, including $2$.

\subsection*{Lean formalization}
The statement was formalized by math0 in the file
\path{lean/Math0/Statements/OEIS/A335012.lean}, which was reviewed for fidelity and frozen (SHA-256 \texttt{954d748d\ldots dbf31}) before any proof attempt.
Its declarations are, verbatim except that the copyright header, the \texttt{import} line, docstrings and blank lines are omitted and the proof of \decl{isLongPeriodPrimeB_iff} is replaced by \texttt{...}:
\begin{lstlisting}[basicstyle=\ttfamily\scriptsize]
namespace Math0.Statements.A335012
def a (n : ℕ) : ℕ := ((Finset.Icc 1 n).image (fun k => ((10 ^ k - 1) / 9) % n)).card
theorem a_data : (List.range 81).map (fun i => a (i + 1)) =
    [1, 1, 3, 2, 1, 3, 6, 3, 9, 1, 2, 4, 6, 6, 3, 4, 16, 9, 18, 2, 6, 2, 22, 5, 2, 6, 27, 7, 28, 3,
      15, 5, 6, 16, 6, 10, 3, 18, 6, 3, 5, 6, 21, 3, 9, 22, 46, 6, 42, 2, 48, 7, 13, 27, 2, 8, 18,
      28, 58, 4, 60, 15, 18, 6, 6, 6, 33, 17, 66, 6, 35, 11, 8, 3, 4, 19, 6, 6, 13, 4, 81] := by
  decide +kernel
def IsLongPeriodPrime (p : ℕ) : Prop :=
  p.Prime ∧ (p = 2 ∨ orderOf (10 : ZMod p) = p - 1)
def isLongPeriodPrimeB (p : ℕ) : Bool :=
  decide p.Prime && (p == 2 || (10 ^ (p - 1) % p == 1 % p &&
    (List.range' 1 (p - 2)).all fun m => 10 ^ m % p != 1 % p))
theorem isLongPeriodPrimeB_iff (p : ℕ) : isLongPeriodPrimeB p = true ↔ IsLongPeriodPrime p := by
  ...
theorem isLongPeriodPrime_data :
    (List.range 270).filter isLongPeriodPrimeB =
      [2, 7, 17, 19, 23, 29, 47, 59, 61, 97, 109, 113, 131, 149, 167, 179, 181, 193, 223, 229,
        233, 257, 263, 269] := by
  decide +kernel
def statement : Prop :=
  ∀ n : ℕ, 1 ≤ n → (a n = n - 1 ↔ IsLongPeriodPrime n)
end Math0.Statements.A335012
\end{lstlisting}

\subsection*{How the definitions encode the OEIS entry}
\begin{itemize}
\item \emph{The sequence.} In \decl{a}, the term \texttt{(10 \^{} k - 1) / 9} is computed in $\N$ with truncated subtraction and floor division; since $10^k\ge1$ and $9\mid 10^k-1$, it equals $R_k$ exactly (this is the lemma \decl{rep_eq} of the proof file). Thus \texttt{a n} is the number of elements of the image of $k\mapsto R_k \bmod n$ on $\{1,\dots,n\}$, which is the definition of A335012 and agrees with the PARI program \texttt{\#Set(vector(n, k, (10\^{}k-1)/9) \% n)} and the Maple programs of the entry. The theorem \decl{a_data} checks, by evaluation in the Lean kernel, that \texttt{a} reproduces all $81$ listed terms $a(1),\dots,a(81)$.
\item \emph{Long period primes.} \decl{IsLongPeriodPrime} says that $p$ is prime and either $p=2$ or $10$ has order $p-1$ in $\Z/p\Z$. For primes $p\notin\{2,5\}$ this is the period condition, as recalled in the introduction. For $p=5$ we have $10=0$ in $\Z/5\Z$, and Mathlib's \decl{orderOf x} is $0$ by convention when no positive power of \texttt{x} equals $1$, so the condition reads $0=4$ and fails; this correctly excludes $5$. For $p=2$ the same convention gives \texttt{orderOf (10 : ZMod 2) = 0}, which is why the disjunct $p=2$ is needed. For $p=3$ the order is $1\ne2$, so $3$ is excluded. These are the same conventions as the Maple program and the alternative Mathematica program of A006883, which test, for primes $p$, whether $p=2$ or $\ord_p(10)=p-1$. As a check, \decl{isLongPeriodPrime_data} shows by kernel evaluation that the numbers below $270$ satisfying the Boolean test \decl{isLongPeriodPrimeB} are exactly the first $24$ listed terms of A006883, and \decl{isLongPeriodPrimeB_iff} proves that this Boolean test is equivalent to \decl{IsLongPeriodPrime}.
\item \emph{Range.} The statement quantifies over all $n\ge1$, the index range of A335012. The hypothesis $1\le n$ is necessary: \texttt{a 0 = 0 = 0 - 1} in $\N$ while $0$ is not prime, so the statement would be false at $n=0$. It adds no other hypotheses.
\end{itemize}

\section{Proof}\label{sec:proof}

The proof follows the Lean file \path{lean/Math0/Frontier/OEIS/A335012.lean}; with each step we give the name of the corresponding Lean declaration.
Following that file, we define the repunits recursively by
\[
  R_0=0,\qquad R_{k+1}=10R_k+1\qquad(k\ge0)
\]
(\decl{rep}). For an integer $x$ and $n\ge1$ we write $x \bmod n\in\{0,\dots,n-1\}$ for the remainder of $x$ upon division by $n$.

\begin{lemma}[\decl{nine_rep}, \decl{rep_eq}, \decl{a_eq}, \decl{rep_add}, \decl{rep_mod_three}, \decl{rep_mod_two}, \decl{rep_mod_five}]\label{lem:basic}
For all $k,m\ge0$:
\begin{enumerate}
\item $9R_k+1=10^k$; hence $(10^k-1)/9=R_k$, and $a(n)=\#\{R_k\bmod n: 1\le k\le n\}$ for every $n$;
\item $R_{m+k}=10^kR_m+R_k$;
\item $R_k\equiv k\pmod 3$;
\item if $k\ge1$, then $R_k\equiv1\pmod 2$ and $R_k\equiv 1\pmod 5$.
\end{enumerate}
\end{lemma}

\begin{proof}
(1) By induction on $k$: $9R_{k+1}+1=10(9R_k+1)=10^{k+1}$.
(2) By induction on $k$: the case $k=0$ is $R_m=R_m$, and $R_{m+k+1}=10R_{m+k}+1=10^{k+1}R_m+(10R_k+1)$.
(3) By induction: $R_{k+1}=10R_k+1\equiv R_k+1\pmod 3$.
(4) $R_k=10R_{k-1}+1\equiv1\pmod{10}$.
\end{proof}

\begin{lemma}[\decl{modEq_of_dvd_rep}, \decl{dvd_rep_of_modEq}]\label{lem:div}
Let $n\ge1$ and $m\ge0$. If $n\mid R_m$, then $10^m\equiv1\pmod n$. Conversely, if $\gcd(n,3)=1$ and $10^m\equiv1\pmod n$, then $n\mid R_m$.
\end{lemma}

\begin{proof}
By Lemma~\ref{lem:basic}(1), $10^m-1=9R_m$. If $n\mid R_m$ then $n\mid 10^m-1$. Conversely, if $10^m\equiv1\pmod n$ then $n\mid 9R_m$; since $\gcd(n,3)=1$ we have $\gcd(n,9)=1$, and hence $n\mid R_m$.
\end{proof}

In Lean the order of $10$ is taken in the ring \texttt{ZMod n}; the lemma \decl{zmod_pow_iff} translates $10^m=1$ in $\Z/n\Z$ into $10^m\equiv1\pmod n$.

\begin{lemma}[periodicity; \decl{exists_small}, \decl{a_le_of_dvd}]\label{lem:per}
Let $n,m\ge1$ with $n\mid R_m$. Then for every $k\ge1$ there is $j$ with $1\le j\le m$ and $R_j\equiv R_k\pmod n$. Consequently $a(n)\le m$.
\end{lemma}

\begin{proof}
By strong induction on $k$. If $k\le m$ take $j=k$. If $k>m$, then $1\le k-m<k$, so by induction there is $j\in[1,m]$ with $R_j\equiv R_{k-m}\pmod n$, and by Lemma~\ref{lem:basic}(2), $R_k=10^{k-m}R_m+R_{k-m}\equiv R_{k-m}\equiv R_j\pmod n$.
Hence $\{R_k\bmod n:1\le k\le n\}\subseteq\{R_j\bmod n:1\le j\le m\}$, and the latter set has at most $m$ elements.
\end{proof}

\begin{lemma}[cancellation; \decl{dvd_rep_sub}]\label{lem:cancel}
Let $n\ge1$ with $\gcd(n,10)=1$, and let $0\le i<j$ with $R_i\equiv R_j\pmod n$. Then $n\mid R_{j-i}$.
\end{lemma}

\begin{proof}
By Lemma~\ref{lem:basic}(2), $R_j=10^iR_{j-i}+R_i$, so $n\mid 10^iR_{j-i}$. Since $\gcd(n,10^i)=1$, $n\mid R_{j-i}$.
\end{proof}

\begin{lemma}[pigeonhole; \decl{exists_collision}]\label{lem:pigeon}
If $a(n)<n$, there are $1\le i<j\le n$ with $R_i\equiv R_j\pmod n$.
\end{lemma}

\begin{proof}
Otherwise $k\mapsto R_k\bmod n$ is injective on $\{1,\dots,n\}$, and its image has $n$ elements.
\end{proof}

\begin{lemma}[\decl{a_le_sub_two}]\label{lem:miss}
Let $n\ge3$, and suppose that $R_k\bmod n\notin\{0,2\}$ for all $k\ge1$. Then $a(n)\le n-2$.
\end{lemma}

\begin{proof}
All residues $R_k\bmod n$ lie in $\{0,1,\dots,n-1\}\setminus\{0,2\}$, a set with $n-2$ elements since $n\ge3$.
\end{proof}

Finally, $a(1)=1$ and $a(2)=1$ (\decl{a_one}, \decl{a_two}): modulo $1$ the only residue is $0$, and $R_1=1$, $R_2=11$ are both odd. In Lean both values are obtained by \texttt{decide +kernel}, that is, by evaluating the definition of \decl{a} in the kernel.

\begin{proposition}[\decl{forward}]\label{prop:forward}
Let $n\ge1$ with $a(n)=n-1$. Then $n$ is a long period prime.
\end{proposition}

\begin{proof}
We consider five cases, in the order of the Lean proof.

\emph{Case $n=1$.} Then $a(1)=1\ne0=n-1$, so this case does not occur.

\emph{Case $2\mid n$.} If $n=2$, then $n$ is a long period prime by definition. Otherwise $n\ge4$. Since $n$ is even, $R_k\bmod n\equiv R_k\equiv1\pmod 2$ for $k\ge1$ by Lemma~\ref{lem:basic}(4), so every residue $R_k\bmod n$ is odd and in particular differs from $0$ and $2$. Lemma~\ref{lem:miss} gives $a(n)\le n-2<n-1$, a contradiction.

\emph{Case $5\mid n$.} Then $n\ge5$, and $R_k\bmod n\equiv R_k\equiv1\pmod5$ for $k\ge1$, so again no residue equals $0$ or $2$, and Lemma~\ref{lem:miss} gives a contradiction.

From now on $n\ge3$ and $\gcd(n,10)=1$.

\emph{Case $3\mid n$.} Since $a(n)=n-1<n$, Lemma~\ref{lem:pigeon} gives $1\le i<j\le n$ with $R_i\equiv R_j\pmod n$, and Lemma~\ref{lem:cancel} gives $n\mid R_m$ with $m=j-i\in[1,n-1]$. By Lemma~\ref{lem:per}, $n-1=a(n)\le m$, so $m=n-1$ and $n\mid R_{n-1}$. Then $3\mid R_{n-1}$, and $R_{n-1}\equiv n-1\pmod 3$ by Lemma~\ref{lem:basic}(3), so $3\mid n-1$. Together with $3\mid n$ this gives $3\mid1$, a contradiction.

\emph{Case $\gcd(n,30)=1$.} By Euler's theorem $10^{\varphi(n)}\equiv1\pmod n$, and $\varphi(n)\ge1$. Lemma~\ref{lem:div} (which applies since $\gcd(n,3)=1$) gives $n\mid R_{\varphi(n)}$, and Lemma~\ref{lem:per} gives $n-1=a(n)\le\varphi(n)$. Since $n\ge2$ we also have $\varphi(n)<n$, so $\varphi(n)=n-1$ and $n$ is prime. Let $d$ be the order of $10$ in $\Z/n\Z$. As $10^{n-1}\equiv1\pmod n$, $d$ divides $n-1$, so $1\le d\le n-1$. Since $10^d\equiv1\pmod n$, Lemma~\ref{lem:div} gives $n\mid R_d$, and Lemma~\ref{lem:per} gives $n-1=a(n)\le d$. Hence $d=n-1$, and $n$ is a long period prime.
\end{proof}

\begin{proposition}[\decl{backward}]\label{prop:backward}
If $p$ is a long period prime, then $a(p)=p-1$.
\end{proposition}

\begin{proof}
If $p=2$, then $a(2)=1=p-1$. Otherwise $p\ge3$ is prime and $10$ has order $d=p-1$ in $\Z/p\Z$; in particular $10^{p-1}\equiv1\pmod p$.

\emph{$p$ is coprime to $30$.} From $10^{p-1}\equiv1\pmod p$ we get $\gcd(10^{p-1},p)=\gcd(1,p)=1$, and since $p-1\ge1$ this gives $\gcd(10,p)=1$. Moreover $p\ne3$: in $\Z/3\Z$ we have $10^1=1$ with $1<2=d$, which contradicts the minimality of the order. Since $p$ and $3$ are distinct primes, $\gcd(p,3)=1$.

\emph{Upper bound.} By Lemma~\ref{lem:div}, $p\mid R_{p-1}$, so $a(p)\le p-1$ by Lemma~\ref{lem:per}.

\emph{Lower bound.} The residues $R_1\bmod p,\dots,R_{p-1}\bmod p$ are pairwise distinct. Indeed, if $R_i\equiv R_j\pmod p$ with $1\le i<j\le p-1$, then $p\mid R_{j-i}$ by Lemma~\ref{lem:cancel}, so $10^{j-i}\equiv1\pmod p$ by Lemma~\ref{lem:div}. But $0<j-i<p-1=d$, which contradicts the minimality of $d$. These $p-1$ distinct residues are among the residues $R_k\bmod p$, $1\le k\le p$, so $a(p)\ge p-1$.
\end{proof}

\begin{proof}[Proof of Theorem~\ref{thm:main} (\decl{result})]
For $n\ge1$, Proposition~\ref{prop:forward} gives the forward implication and Proposition~\ref{prop:backward} the converse.
\end{proof}

\begin{remark}[Israel's formula]\label{rem:israel}
The same lemmas give the formula $a(n)=\ord_n(10)$ for $\gcd(n,30)=1$ stated by Israel in A335012.
Let $d=\ord_n(10)$, which exists since $\gcd(n,10)=1$. By Lemma~\ref{lem:div}, $n\mid R_d$, so $a(n)\le d$ by Lemma~\ref{lem:per}. As in the lower bound of Proposition~\ref{prop:backward}, Lemmas~\ref{lem:cancel} and~\ref{lem:div} and the minimality of $d$ show that $R_1,\dots,R_d$ are pairwise distinct modulo $n$; since $d\mid\varphi(n)$ we have $d\le n$, so $a(n)\ge d$.
This remark is not part of the Lean development, which avoids the order of $10$ modulo a composite $n$ by using Euler's theorem directly.
\end{remark}

\section{Machine verification}\label{sec:lean}

\subsection*{Files and declarations}
The proof is in the file \path{lean/Math0/Frontier/OEIS/A335012.lean} of the math0 repository, \url{https://github.com/Foundation-0/math0}. It imports the frozen statement file \path{lean/Math0/Statements/OEIS/A335012.lean} (Section~\ref{sec:statement}) and the audit commands in \path{lean/Math0/Audit.lean}. The main declaration is \decl{Math0.Frontier.A335012.result}, of type \decl{Math0.Statements.A335012.statement}. All auxiliary declarations, named in Section~\ref{sec:proof}, are private. The repunit, the two directions and the end of the file are, verbatim (docstrings, blank lines and the proofs of the two directions omitted):
\begin{lstlisting}
private def rep : ℕ → ℕ
  | 0 => 0
  | k + 1 => 10 * rep k + 1
private lemma forward {n : ℕ} (hn : 1 ≤ n) (ha : a n = n - 1) : IsLongPeriodPrime n := by
  ...
private lemma backward {p : ℕ} (hp : IsLongPeriodPrime p) : a p = p - 1 := by
  ...
theorem result : statement := fun _ hn => ⟨forward hn, backward⟩
end Math0.Frontier.A335012
#math0_settles Math0.Frontier.A335012.result Math0.Statements.A335012.statement
\end{lstlisting}
The proof file uses no \texttt{sorry}, \texttt{native\_decide} or \texttt{set\_option}. Its only proofs by evaluation are \texttt{decide +kernel} for $a(1)=1$ and $a(2)=1$, and one \texttt{decide} for $10^1=1$ in \texttt{ZMod 3}; both tactics produce proofs checked by the kernel and add no axioms. Euler's theorem (\decl{Nat.ModEq.pow_totient}), the fact that $\varphi(n)=n-1$ if and only if $n$ is prime, for $n\ge1$ (\decl{Nat.totient_eq_iff_prime}), and the basic properties of \decl{orderOf} come from Mathlib.

\subsection*{Audit}
The command \texttt{\#math0\_settles}, defined in \path{lean/Math0/Audit.lean} for statements that math0 formalizes itself, fails unless (1) \decl{statement} is a \texttt{Prop}-valued definition without universe parameters; (2) \decl{statement} and \decl{result} use no axioms besides \decl{propext}, \decl{Classical.choice} and \decl{Quot.sound} (this rejects \decl{sorryAx}, \decl{Lean.ofReduceBool} from \texttt{native\_decide}, and custom axioms); and (3) the type of \decl{result} is syntactically the statement or its negation, up to metadata, with no unfolding. Here it reports
\begin{lstlisting}[basicstyle=\ttfamily\scriptsize]
math0 audit: 'Math0.Statements.A335012.statement' depends on axioms [Classical.choice, Quot.sound, propext]
math0 audit: 'Math0.Frontier.A335012.result' depends on axioms [Classical.choice, Quot.sound, propext]
math0 statement check: 'Math0.Frontier.A335012.result' proves 'Math0.Statements.A335012.statement'
\end{lstlisting}

\subsection*{Toolchain and reproduction}
We use Lean~4 \cite{Lean4} \texttt{v4.33.1} and Mathlib \cite{Mathlib} \texttt{v4.33.1}, pinned in \path{lean/}; the repository also pins formal-conjectures \cite{FC} at commit \texttt{c77fdd4}, which this result does not use.
To reproduce, run \path{scripts/verify.sh} from the repository root: it fetches the Mathlib build cache, builds all math0 results with \texttt{lake build} (re-running every audit command), checks that the pinned \texttt{formal\_conjectures} and \texttt{mathlib} checkouts are unmodified, and reports the number of passed statement checks. Once its imports are built, the single file can be checked with \texttt{lake env lean Math0/Frontier/OEIS/A335012.lean} in the directory \path{lean}.

{\small
\begin{thebibliography}{99}

\bibitem{FC}
M.~Firsching, P.~Lezeau, S.~Mercuri, M.~Z. Horv\'ath, Y.~Dillies, C.~S\"onne, E.~Wieser, F.~Zhang, T.~Hubert, B.~Ag\"uera y Arcas, and P.~Kohli, \emph{Formal Conjectures: An Open and Evolving Benchmark for Verified Discovery in Mathematics}, arXiv:2605.13171, 2026. Repository: \url{https://github.com/google-deepmind/formal-conjectures}.

\bibitem{HW}
G.~H. Hardy and E.~M. Wright, \emph{An Introduction to the Theory of Numbers}, 6th ed., revised by D.~R. Heath-Brown and J.~H. Silverman, Oxford University Press, 2008.

\bibitem{Mathlib}
The mathlib Community, \emph{The Lean mathematical library}, in: Proceedings of the 9th ACM SIGPLAN International Conference on Certified Programs and Proofs (CPP 2020), ACM, 2020, 367--381.

\bibitem{Lean4}
L.~de Moura and S.~Ullrich, \emph{The Lean 4 theorem prover and programming language}, in: Automated Deduction -- CADE 28, Lecture Notes in Comput. Sci. \textbf{12699}, Springer, 2021, 625--635.

\bibitem{MW}
E.~W. Weisstein, \emph{Full Reptend Prime}, MathWorld -- A Wolfram Web Resource, \url{https://mathworld.wolfram.com/FullReptendPrime.html}.

\bibitem{A001913}
The OEIS Foundation Inc., \url{https://oeis.org/A001913} (full reptend primes: primes with primitive root 10), entry by N.~J.~A. Sloane.

\bibitem{A002275}
The OEIS Foundation Inc., \url{https://oeis.org/A002275} (repunits), entry by N.~J.~A. Sloane.

\bibitem{A006883}
The OEIS Foundation Inc., \url{https://oeis.org/A006883} (long period primes), entry by R.~Munafo.

\bibitem{A084680}
The OEIS Foundation Inc., \url{https://oeis.org/A084680} (order of 10 modulo $n$), entry by L.~Beedassy (2003).

\bibitem{A335012}
The OEIS Foundation Inc., \url{https://oeis.org/A335012}, entry by S.-P. Eu (May 19, 2020); comment by R.~Israel (June 25, 2020); revision \#31, Mar.~12, 2026.

\end{thebibliography}
}
\end{document}
